\section{Excluding flat faces}
\label{conclusion:section}

\begin{proof}[Proof of Theorem~\ref{thm:exponential-shape}]
Suppose $\partial\Bcal$ contains a nontrivial segment, and fix the
maximal supporting face with extreme endpoints $a,b$ described in
Section~\ref{model:section}. The same admissible road class has the
bounds
\[
 0<c_0\le b(n)\le C_s n^{3/4}
\]
from Propositions~\ref{prop:fb-positive} and \ref{prop:fb-upper}.

Proposition~\ref{improvement:contraction-proposition} supplies fixed
rational directions, $s>0$, an integer $D\ge1$, and $\rho<1$ such that
$b(Dn)\le\rho b(n)$ for every sufficiently large admissible $n$.
Choose an admissible $n_0$ above the common threshold. Multiplication
by $D$ preserves admissibility, so iteration gives
\[
 0<c_0\le b(D^k n_0)\le\rho^k b(n_0)
 \qquad(k\ge0),
\]
a contradiction. Thus every supporting face is a single point.

To express this as strict convexity, take distinct
$x,y\in\partial\Bcal$ and $0<t<1$. Convexity gives
$\mu((1-t)x+ty)\le1$. Equality would provide a supporting functional
$\ell\le\mu$ at this point, and therefore
\[
 1=(1-t)\ell(x)+t\ell(y),\qquad \ell(x),\ell(y)\le1.
\]
Both endpoint values would be one, so the entire segment $[x,y]$
would lie in a nontrivial supporting face. This has been excluded.
Therefore $\mu((1-t)x+ty)<1$ for the rate-one exponential law.

Theorem~\ref{thm:gamma-differentiability} with shape $\kappa=1$
gives Fr\'echet differentiability of this norm away from the origin and
a $C^1$ unit sphere, with a unique supporting line at every point.
To pass to any rate $\lambda>0$, couple the weights by
$\tau_e^{(\lambda)}=\lambda^{-1}\tau_e^{(1)}$. Then every path cost,
and hence every passage time, is multiplied by $\lambda^{-1}$. Thus
\[
 T_\lambda(x,y)=\lambda^{-1}T_1(x,y),\qquad
 \mu_\lambda=\lambda^{-1}\mu_1,\qquad
 \Bcal_\lambda=\lambda\Bcal_1.
\]
Multiplication by a positive scalar preserves differentiability,
strict convexity, and the $C^1$ boundary property. This proves every
assertion of the theorem.
\end{proof}
